\documentclass[11pt,a4paper]{article}
\usepackage[utf8]{inputenc}
\usepackage[T1]{fontenc}
\usepackage{amsmath, amssymb}
\usepackage{geometry}
\usepackage{enumitem}
\usepackage{hyperref}

\geometry{margin=1in}

\title{Simple Harmonic Oscillator and Harmonic Regression for NDVI}
\author{}
\date{}

\begin{document}

\maketitle

The general formula for the displacement $x(t)$ of a \textbf{simple harmonic oscillator} over time is:
\begin{equation}
x(t) = A \cos(\omega t + \phi)
\end{equation}

\noindent \textit{Note: This formulation is applied directly in remote sensing as \textbf{harmonic regression} to fit smooth seasonal sine and cosine curves to noisy NDVI time-series data.}

\section*{1. Key Variables in the Formula}
\begin{itemize}[leftmargin=1.5em]
    \item $x(t)$: Position (or value, such as NDVI) at time $t$.
    \item $A$: \textbf{Amplitude} --- the maximum displacement from the equilibrium (center) position.
    \item $\omega$: \textbf{Angular frequency} --- the rate of oscillation, measured in radians per unit time ($\omega = 2\pi f = \frac{2\pi}{T}$, where $f$ is frequency and $T$ is the period).
    \item $t$: \textbf{Time}.
    \item $\phi$: \textbf{Phase constant (phase shift)} --- determines the starting point of the wave at $t = 0$.
\end{itemize}

\section*{2. Core Governing Equations (Physics)}
In classical mechanics, a simple harmonic oscillator is driven by a restoring force proportional to displacement (\textbf{Hooke's Law}):
\begin{itemize}[leftmargin=1.5em]
    \item \textbf{Restoring Force:}
    \begin{equation}
    F = -kx
    \end{equation}
    \item \textbf{Differential Equation of Motion:}
    \begin{equation}
    \frac{d^2x}{dt^2} + \omega^2 x = 0
    \end{equation}
    \item \textbf{Angular Frequency (Mass-Spring System):}
    \begin{equation}
    \omega = \sqrt{\frac{k}{m}}
    \end{equation}
    where $k$ is the spring constant and $m$ is mass.
\end{itemize}

\section*{3. Harmonic Form Used for NDVI Time Series}
When modeling annual vegetation curves (harmonic analysis / Fourier series), the formula is typically expressed as a linear combination of sine and cosine waves plus a baseline offset ($c_0$):
\begin{equation}
\mathrm{NDVI}(t) = c_0 + \sum_{n=1}^{N} \left[ a_n \cos\left(\frac{2\pi n t}{T}\right) + b_n \sin\left(\frac{2\pi n t}{T}\right) \right]
\end{equation}

\begin{itemize}[leftmargin=1.5em]
    \item $c_0$: Mean annual NDVI baseline.
    \item $T$: Period (e.g., $365.25$ days for annual cycles).
    \item $n = 1$: Captures single-peak annual seasons (e.g., deciduous forests, single-crop fields).
    \item $n = 2$: Captures bimodal seasons (e.g., double-cropping or bimodal rainfall regimes).
\end{itemize}

\end{document}
